\scanpage{79}
$\mathbb P(\widetilde{\mathfrak g}/\mathcal O_S)$, \deleted{\unsure{et}} d'où
\[
X\longrightarrow\mathbb P(\widetilde{\mathfrak g}/\mathcal O_S)\quad
\begin{aligned}&\text{qui est lui-$=$ un}\\&\text{fibré en droites projectives}\end{aligned}\tag*{(\unsure{69|60|59})}
\]
Si \inserted{(59) donc} ce morphisme \inserted{(60)} était encore un monomorphisme, il \deleted{aurait} résulterait que (60) est $=$ un isomorphisme, donc $X$ serait ``banal'' du pt de vue Brauer -- or il n'est pas vrai qu'en car $2$ les fibrés en droites projectives soient tjrs banals !

En fait, on trouve ici un iso. canonique
\[
\mathbb P(\widetilde{\mathfrak g}/\mathcal O_S)\simeq X^{(2)}\;\raisebox{.65ex}{\(\scriptstyle\longrightarrow\text{(frobénisé)}\)}.\tag*{(61)}
\]
et (60) n'est autre que le morphisme canonique
\[
X\longrightarrow X^{(2)}
\]
qui est un épimorphisme radiciel (et \deleted{plat} \inserted{fini loc. libre} de rang $2^d=2$), mais en effet nullement un monomorphisme \ldots.
% archive sheet number: 78
